

\problem{1}
На кошачьей выставке каждый посетитель погладил ровно трёх кошек. При этом оказалось, что каждую кошку погладили ровно три посетителя. Докажите, что посетителей было ровно столько же, сколько кошек.

\problem{2}
В 7 классе каждый мальчик дружит ровно с двумя девочками, а каждая девочка - ровно с тремя мальчиками. Ещё известно, что в классе 19 парт, а вчера 31 семиклассник побывал в Палеонтологическом музее. Сколько человек в этом классе?

\problem{3}
Можно ли в клетки квадрата $10 \times 10$ поставить некоторое количество звёздочек так, чтобы в каждом квадрате $2 \times 2$ было ровно две звёздочки, а в каждом прямоугольнике $3 \times 1$ - ровно одна звёздочка?

\problem{4}
Аркаша вырезал много одинаковых квадратов и в вершинах каждого из них в произвольном порядке написал числа $1,2,3$ и $4.3 а т е м$ Степан сложил квадраты в стопку и написал сумму чисел, попавших в каждый из четырех углов стопки. Может ли оказаться так, что в каждом углу стопки сумма равна
a) 2015?

6) 2016 ?

\problem{5}
Можно ли в прямоугольной таблице $5 \times 10$ так расставить числа, чтобы сумма чисел каждой строки равнялась бы 30 , а сумма чисел каждого столбца равнялась бы 10 ?

\problem{6}
а) Можно ли занумеровать рёбра куба числами $-6,-5,-4,-3,-2,-1,1,2,3,4,5,6$ так, чтобы для каждой вершины куба сумма номеров рёбер, которые в ней сходятся, была одинаковой?

6) А можно ли так же занумеровать рёбра куба натуральными числами от 1 до 12?

\problem{7}
Представьте число 2015 как разность двух палиндромов. (Палиндромы - это числа, не меняющиеся при прочтении справа налево, например, $0,717,2002$.)